% from main.tex
%%%%%%%%%%  scan idx 623  |  folio 610  %%%%%%%%%%
Pour \(i,j\ge0\), \(R^if_*\mathcal O_X\) et \(R^jg_*\mathcal O_{X\times_YX}\) sont plats sur \(S\).
\item La suite
\begin{equation}\tag{*}
  \mathcal O_Y\longrightarrow f_*\mathcal O_X
  \mathrel{\substack{\xrightarrow{\pi_1}\\[-0.7ex]\xrightarrow[\pi_2]{}}}
  g_*\mathcal O_{X\times_YX}
\end{equation}
est exacte et \(\coker(\pi_1-\pi_2)\) est un \(\mathcal O_Y\)-module plat sur \(S\).
\end{enumerate}
Alors \(f^*:\Pic_{Y/S}\longrightarrow\Pic_{X/S}\) est représentable par un morphisme quasi-affine de présentation finie.

En effet notons d'abord que 1)-4) sont stables par tout changement de base \(T\to S\). Seule la stabilité de 3) et 4) n'est pas évidente. Mais d'après la formule de Künneth (EGA III 6.9.9), 1) et 3) impliquent que la formation des \(R^if_*(\mathcal O_X)\) et \(R^ig_*(\mathcal O_{X\times_YX})\) commutent aux changements de base \(T\to S\), ce qui implique aussitôt la stabilité des conditions 3). De plus, la formation du conoyau de \(\pi_1-\pi_2:f_*\mathcal O_X\longrightarrow g_*\mathcal O_{X\times_YX}\) commute donc aussi aux changements de base~; comme il est plat sur \(S\) (d'après 4)), \(\im(\pi_1-\pi_2)\) est plat sur \(S\) et commute aux changements de base. Par suite, \(\Ker(\pi_1-\pi_2)\) commute aux changements de base \(T\to S\), et est donc égal à \(\mathcal O_Y\) universellement.

La stabilité des hypothèses implique deux choses:a. L'assertion à démontrer étant locale pour f.p.p.f., on peut par descente supposer que \(X\) possède une section \(\alpha\), donc \(Y\) possède la section \(\beta=f\circ\alpha\). b. Pour établir l'assertion il suffit de considérer un \(S\)-point de \(\Pic_{X/S}\), c'est-à-dire d'après a) et 2) un faisceau inversible \(L\) sur \(X\) muni d'une rigidification \(u\) relativement à \(\alpha\), et de montrer que le faisceau f.p.p.f. \(Z\) image réciproque de ce point dans \(\Pic_{Y/S}\) est représentable par un \(S\)-schéma quasi-affine, de présentation finie sur \(S\).

